% Emacs settings: -*-mode: latex; TeX-master: "manual.tex"; -*-

\chapter{Monochromators}

In this class of components, we are concerned with elastic Bragg
scattering from monochromators. \textbf{Monochromator\_flat}
models a flat thin mosaic crystal with a single scattering vector
perpendicular to the surface.
The component \textbf{Monochromator\_curved} is physically similar,
but models a singly or doubly bend monochromator crystal arrangement.

A much more general model of scattering from a single crystal is
found in the component \textbf{Single\_crystal},
which is presented under Samples, chapter~\ref{c:samples}.

\section{The \texttt{Monochromator\_flat} McStas Component}
Flat Monochromator crystal with anisotropic mosaic.

\subsection*{Identification}
\begin{itemize}
  \item \textbf{Site:} 
  \item \textbf{Author:} Kristian Nielsen
  \item \textbf{Origin:} Risoe
  \item \textbf{Date:} 1999
\end{itemize}

\subsection*{Description}
\begin{lstlisting}
Flat, infinitely thin mosaic crystal, useful as a monochromator or analyzer.
For an unrotated monochromator component, the crystal surface lies in the Y-Z
plane (ie. parallel to the beam).
The mosaic is anisotropic gaussian, with different FWHMs in the Y and Z
directions. The scattering vector is perpendicular to the surface.

Example: Monochromator_flat(zmin=-0.1, zmax=0.1, ymin=-0.1, ymax=0.1, mosaich=30.0, mosaicv=30.0, r0=0.7, Q=1.8734)

Monochromator lattice parameter
PG       002 DM=3.355 AA (Highly Oriented Pyrolythic Graphite)
PG       004 DM=1.677 AA
Heusler  111 DM=3.362 AA (Cu2MnAl)
CoFe         DM=1.771 AA (Co0.92Fe0.08)
Ge       111 DM=3.266 AA
Ge       311 DM=1.714 AA
Ge       511 DM=1.089 AA
Ge       533 DM=0.863 AA
Si       111 DM=3.135 AA
Cu       111 DM=2.087 AA
Cu       002 DM=1.807 AA
Cu       220 DM=1.278 AA
Cu       111 DM=2.095 AA
\end{lstlisting}

\subsection*{Input parameters}
Parameters in \textbf{boldface} are required; the others are optional.

\begin{longtable}{p{0.22\textwidth}p{0.12\textwidth}p{0.46\textwidth}p{0.14\textwidth}}
\toprule
\textbf{Name} & \textbf{Unit} & \textbf{Description} & \textbf{Default} \\
\midrule
\endhead
zmin & m & Lower horizontal (z) bound of crystal & -0.05 \\
zmax & m & Upper horizontal (z) bound of crystal & 0.05 \\
ymin & m & Lower vertical (y) bound of crystal & -0.05 \\
ymax & m & Upper vertical (y) bound of crystal & 0.05 \\
zwidth & m & Width of crystal, instead of zmin and zmax & 0 \\
yheight & m & Height of crystal, instead of ymin and ymax & 0 \\
mosaich & arc minutes & Horizontal mosaic (in z direction) (FWHM) & 30.0 \\
mosaicv & arc minutes & Vertical mosaic (in y direction) (FWHM) & 30.0 \\
r0 & 1 & Maximum reflectivity & 0.7 \\
Q & 1/angstrom & Magnitude of scattering vector & 1.8734 \\
DM & AA & monochromator d-spacing, instead of Q = 2*pi/DM & 0 \\
\bottomrule
\end{longtable}

\subsection*{Links}
\begin{itemize}
  \item \href{run:/home/willend/willend-McCode/mcstas-comps/optics/Monochromator_flat.comp}{Source code} for \texttt{Monochromator\_flat.comp}.
\end{itemize}
\IfFileExists{Monochromator_flat_static.tex}{\input{Monochromator_flat_static.tex}}{}

\section{The \texttt{Monochromator\_curved} McStas Component}
Double bent multiple crystal slabs with anisotropic gaussian mosaic.

\subsection*{Identification}
\begin{itemize}
  \item \textbf{Site:} 
  \item \textbf{Author:} Emmanuel Farhi, Kim, Lefmann, Peter Link
  \item \textbf{Origin:} \textless{}a href="http://www.ill.fr"\textgreater{}ILL\textless{}/a\textgreater{}
  \item \textbf{Date:} Aug. 24th 2001
\end{itemize}

\subsection*{Description}
\begin{lstlisting}
Double bent infinitely thin mosaic crystal, useful as a monochromator or
analyzer. which uses a small-mosaicity approximation and taking into account
higher  order scattering. The mosaic is anisotropic gaussian, with different
FWHMs in the Y and Z directions. The scattering vector is perpendicular to the
surface. For an unrotated monochromator component, the crystal plane lies in
the y-z plane (ie. parallel to the beam). The component works in reflection, but
also transmits the non-diffracted beam. Reflectivity and transmission files may
be used. The slabs are positioned in the vertical plane (not on a
cylinder/sphere), and are rotated according to the curvature radius.
When curvatures are set to 0, the monochromator is flat.
The curvatures approximation for parallel beam focusing to distance L, with
monochromator rotation angle A1 are:
RV = 2*L*sin(DEG2RAD*A1);
RH = 2*L/sin(DEG2RAD*A1);

When you rotate the component by A1 = asin(Q/2/Ki)*RAD2DEG, do not forget to
rotate the following components by A2=2*A1 (for 1st order) !

Example: Monochromator_curved(zwidth=0.01, yheight=0.01, gap=0.0005, NH=11, NV=11, mosaich=30.0, mosaicv=30.0, r0=0.7, Q=1.8734)

Monochromator lattice parameter
PG       002 DM=3.355 AA (Highly Oriented Pyrolythic Graphite)
PG       004 DM=1.677 AA
Heusler  111 DM=3.362 AA (Cu2MnAl)
CoFe         DM=1.771 AA (Co0.92Fe0.08)
Ge       111 DM=3.266 AA
Ge       311 DM=1.714 AA
Ge       511 DM=1.089 AA
Ge       533 DM=0.863 AA
Si       111 DM=3.135 AA
Cu       111 DM=2.087 AA
Cu       002 DM=1.807 AA
Cu       220 DM=1.278 AA
Cu       111 DM=2.095 AA
\end{lstlisting}

\subsection*{Input parameters}
Parameters in \textbf{boldface} are required; the others are optional.

\begin{longtable}{p{0.22\textwidth}p{0.12\textwidth}p{0.46\textwidth}p{0.14\textwidth}}
\toprule
\textbf{Name} & \textbf{Unit} & \textbf{Description} & \textbf{Default} \\
\midrule
\endhead
reflect & str & reflectivity file name of text file as 2 columns [k, R] & "NULL" \\
transmit & str & transmission file name of text file as 2 columns [k, T] & "NULL" \\
zwidth & m & horizontal width of an individual slab & 0.01 \\
yheight & m & vertical height of an individual slab & 0.01 \\
gap & m & typical gap  between adjacent slabs & 0.0005 \\
NH & int & number of slabs horizontal & 11 \\
NV & int & number of slabs vertical & 11 \\
mosaich & arc minutes & Horisontal mosaic FWHM & 30.0 \\
mosaicv & arc minutes & Vertical mosaic FWHM & 30.0 \\
r0 & 1 & Maximum reflectivity. O unactivates component & 0.7 \\
t0 & 1 & transmission efficiency & 1.0 \\
Q & AA-1 & Scattering vector & 1.8734 \\
RV & m & radius of vertical focussing. flat for 0 & 0 \\
RH & m & radius of horizontal focussing. flat for 0 & 0 \\
DM & AA & monochromator d-spacing instead of Q=2*pi/DM & 0 \\
mosaic & arc minutes & sets mosaich=mosaicv & 0 \\
width & m & total width of monochromator, along Z & 0 \\
height & m & total height of monochromator, along Y & 0 \\
verbose & 0/1 & verbosity level & 0 \\
order & 1 & specify the diffraction order, 1 is usually prefered. Use 0 for all & 0 \\
\bottomrule
\end{longtable}

\subsection*{Links}
\begin{itemize}
  \item \href{run:/home/willend/willend-McCode/mcstas-comps/optics/Monochromator_curved.comp}{Source code} for \texttt{Monochromator\_curved.comp}.
  \item \textless{}a href="http://mcstas.risoe.dk/pipermail/neutron-mc/1999q1/000133.html"\textgreater{}Additional note\textless{}/a\textgreater{} from Peter Link.
  \item Obsolete Mosaic\_anisotropic by Kristian Nielsen
  \item Contributed Monochromator\_2foc by Peter Link
\end{itemize}
\IfFileExists{Monochromator_curved_static.tex}{\input{Monochromator_curved_static.tex}}{}

\section{Monochromator\_pol: A polarizing monochromator}
\index{Optics!Monochromator, polarizing}
\index{Polarisation!Monochromator\_pol}

The \textbf{Monochromator\_pol} component is the polarising analogue of
\textbf{Monochromator\_flat} above -- a flat, thin mosaic crystal with
spin-dependent (up/down) reflectivity. Its underlying physics (the
nuclear/magnetic structure factor formalism, $F_N$/$F_M$, and how they
relate to the up/down reflectivities) is described in the User Manual's
Polarisation chapter, section on Polarising Monochromator and Guides
(\ref{sub:mono}), which also documents \textbf{Pol\_mirror},
\textbf{Pol\_bender}, \textbf{Pol\_guide\_mirror} and
\textbf{Pol\_guide\_vmirror} -- see also section~\ref{sec:polarising-optics}
of this manual for the rest of the polarisation-aware component set.

\section{The \texttt{Monochromator\_pol} McStas Component}
Flat polarizaing monochromator crystal.

\subsection*{Identification}
\begin{itemize}
  \item \textbf{Site:} 
  \item \textbf{Author:} Peter Christiansen
  \item \textbf{Origin:} RISOE
  \item \textbf{Date:} 2006
\end{itemize}

\subsection*{Description}
\begin{lstlisting}
Based on Monochromator_flat.
Flat, infinitely thin mosaic crystal, useful as a monochromator or analyzer.
For an unrotated monochromator component, the crystal surface lies in the Y-Z
plane (ie. parallel to the beam).
The mosaic and d-spread distributions are both Gaussian.
Neutrons are just reflected (billard ball like). No correction is done for
mosaicity of reflecting crystal.
The crystal is assumed to be a ferromagnet with spin pointing up
eta-tilde = (0, 1, 0) (along y-axis), so that the magnetic field is
pointing opposite (0, -|B|, 0).

The polarisation is done by defining the reflectivity for spin up
(Rup) and spin down (Rdown) (which can be negative, see now!) and
based on this the nuclear and magnetic structure factors are
calculated:
FM = sign(Rup)*sqrt(|Rup|) + sign(Rdown)*sqrt(|Rdown|)
FN = sign(Rup)*sqrt(|Rup|) - sign(Rdown)*sqrt(|Rdown|)
and the physics is calculated as
Pol in = (sx_in, sy_in, sz_in)
Reflectivity R0 = FN*FN + 2*FN*FM*sy_in + FM*FM
(= |Rup| + |Rdown| (for sy_in=0))
Pol out:
sx = (FN*FN - FM*FM)*sx_in/R0;
sy = ((FN*FN - FM*FM)*sy_in + 2*FN*FM + FM*FM*sy_in)/R0;
sz = (FN*FN - FM*FM)*sz_in/R0;

These equations are taken from:
Lovesey: "Theory of neutron scattering from condensed matter, Volume
2", Eq. 10.96 and Eq. 10.110

This component works with gravity (uses PROP_X0).

Example: Monochromator_pol(zwidth=0.2, yheight=0.2, mosaic=30.0, dspread=0.0025, Rup=1.0, Rdown=0.0, Q=1.8734)

Monochromator lattice parameter
PG       002 DM=3.355 AA (Highly Oriented Pyrolythic Graphite)
PG       004 DM=1.677 AA
Heusler  111 DM=3.362 AA (Cu2MnAl)
CoFe         DM=1.771 AA (Co0.92Fe0.08)
Ge       111 DM=3.266 AA
Ge       311 DM=1.714 AA
Ge       511 DM=1.089 AA
Ge       533 DM=0.863 AA
Si       111 DM=3.135 AA
Cu       111 DM=2.087 AA
Cu       002 DM=1.807 AA
Cu       220 DM=1.278 AA
Cu       111 DM=2.095 AA
\end{lstlisting}

\subsection*{Input parameters}
Parameters in \textbf{boldface} are required; the others are optional.

\begin{longtable}{p{0.22\textwidth}p{0.12\textwidth}p{0.46\textwidth}p{0.14\textwidth}}
\toprule
\textbf{Name} & \textbf{Unit} & \textbf{Description} & \textbf{Default} \\
\midrule
\endhead
zwidth & m & Width of crystal & 0.1 \\
yheight & m & Height of crystal & 0.1 \\
mosaic &  & Mosaicity (FWHM) [arc minutes] & 30.0 \\
dspread & 1 & Relative d-spread (FWHM) & 0 \\
Q & AA-1 & Magnitude of scattering vector & 1.8734 \\
DM & AA & monochromator d-spacing instead of Q = 2*pi/DM & 0 \\
pThreshold &  & if probability\textgreater{}pThreshold then accept and weight else random & 0 \\
Rup & 1 & Reflectivity of neutrons with polarization up & 1 \\
Rdown & 1 & Reflectivity of neutrons with polarization down & 1 \\
debug &  & if debug \textgreater{} 0, print out some info about the calculations & 0 \\
\bottomrule
\end{longtable}

\subsection*{Links}
\begin{itemize}
  \item \href{run:/home/willend/willend-McCode/mcstas-comps/optics/Monochromator_pol.comp}{Source code} for \texttt{Monochromator\_pol.comp}.
\end{itemize}
\IfFileExists{Monochromator_pol_static.tex}{\input{Monochromator_pol_static.tex}}{}

\section{Single\_crystal: Thick single crystal monochromator plate with multiple scattering}
\index{Optics!Monochromator, thick}

The \textbf{Single\_crystal} component may be used to study more complex monochromators, including incoherent scattering, thickness and multiple scattering. Please refer to section \ref{s:Single_crystal}.

\section{Phase space transformer - moving monochromator}
\index{Optics!phase space transformer}

Eventhough there exist a few attempts to write dedicated phase space transformer components, there is an elegant way to put a monochromator into move, by mean of the EXTEND keyword. If you define a SPEED parameter for the instrument, the idea is to change the coordinate system before the monochromator, and restore it afterwards, as follow in the TRACE section:
\begin{lstlisting}
DEFINE INSTRUMENT PST(SPEED=200, ...)
(...)
TRACE
(...)
COMPONENT Mono_PST_on=Arm()
AT ...
EXTEND %{
  vx = vx + SPEED; // monochromator moves transversaly by SPPED m/s
%}

COMPONENT  Mono=Monochromator(...)
AT (0,0,0) RELATIVE PREVIOUS

COMPONENT Mono_PST_off=Arm
AT (0,0,0) RELATIVE PREVIOUS
EXTEND  %{
  vz = vz - SPEED; // puts back neutron in static coordinate frame
%}
\end{lstlisting}
This solution does not contain acceleration, but is far enough for most
studies, and it is very simple. In the latter example, the instance \verb+Mono_PST_on+ should itself be rotated to reflect according to a Bragg law.
